% BEGIN THERMAL_R03_SYNC_01
\section{Introduction}

\subsection{Détermination des constantes et relations entre action, aire et information}

Cet article détermine les sections d'action de Planck et leur retour,
la fonctionnelle de Barbero--Immirzi et les coordonnées angulaires CKM, et
établit la transduction de Boltzmann et la réciprocité gravitationnelle dans
les réalisations dimensionnelles spécifiées. Ces résultats sont liés
par des opérations effectives : l'action convertit la fréquence en énergie ;
Barbero--Immirzi pondère l'incidence qui résout l'aire ; Boltzmann
convertit l'information sans dimension en entropie ; et la gravitation
transforme l'action en une échelle géométrique commune. La question physique
est de savoir comment ces fonctions et leurs valeurs sont déterminées à partir d'une même
structure discrète, au lieu d'intervenir comme paramètres mesurés.

La détermination de Planck sépare l'ordre décimal de l'action de ses
coordonnées de retour. Le quotient local du registre possède seize
feuillets ; sa norme multiplicative, suivie de l'auto-échelle, donne
\(\operatorname{ord}_{10}(\varphi\alpha^{16})=-34\).
Les caractères régionaux fixent ensuite
\(\hbar_{\rm pre}=H_5\,10^{-34}\mathcal U_S\) et
\(\hbar_{\rm ret}=\eta_{\rm ret}\,10^{-34}\mathcal U_S\),
avec \(0<\eta_{\rm ret}<H_5\). Le quotient
\(R_{\rm act}=H_5/\eta_{\rm ret}>1\) conserve l'effet du
retour lors du changement de la base positive commune ; dans chaque section,
\(h_a=2\pi\hbar_a\). Le lecteur local, le caractère et la règle de
continuation régionale font partie de cette détermination, développée
en \ref{sec:action}.

La fonctionnelle de Barbero--Immirzi évalue une chaîne orientée de douze
secteurs et une couture de retour sur des incidences de cardinaux 90
et 120. On obtient
\(\gamma_{\HMT}=0.2740076726944592747\ldots\), de parité impaire
sous inversion angulaire. Les règles sectorielles CKM produisent, quant à
elles, une matrice rationnelle de déterminant \(-3857/6\), trois angles
de mélange et la phase \(\delta_{\rm deg}=9A_{\rm deg}\).
Sa réalisation complexe est unitaire et détermine l'invariant de
Jarlskog. Les deux résultats reçoivent les coordonnées angulaires déjà
engendrées : la fonctionnelle de Barbero ne sélectionne pas les angles de mélange,
et ceux-ci ne sélectionnent pas rétrospectivement les règles sectorielles.

La composition longitudinale précède l'évaluation gravitationnelle.
L'inscription des incidences et l'archive complémentaire produisent
\(r_N=5759/23040\). La règle longitudinale R4 applique ce coefficient
à \(\alpha_{\HMT}^{16}L_*\) et détermine \(L\) ; la règle métrique
R5 détermine l'opérateur radial positif du même caractère qu'utilise
l'énergie. La section~\ref{g26:sec:rigidez-radial-es} énonce les deux
compositions et démontre l'unicité de leur réalisation conjointe avec
l'action de phase, l'horloge conjuguée et l'identification du rayon
réduit R8. La référence \(L_*\), déclarée égale à un mètre dans
l'évaluation SI, conserve sa fonction dimensionnelle.

La relation entre Boltzmann et gravitation se concrétise alors dans
\[
 k_B\Theta_{{\rm clk},a}\log3=\frac{h_a}{T_{108}},\qquad
 G_a=\frac{c_{\rm int}^3R_{108}^2}{\hbar_a},\qquad
 \frac{G_a\hbar_a}{c_{\rm int}^3}=\ell_P^2.
\]
Ici, \(R_{108}=L\) et \(t_0=\pi_{\HMT}L/(54c_{\rm int})\).
La première égalité détermine l'énergie par décision tritique et un
transducteur thermique unique une fois fixée la section positive de
température. La deuxième est la spécialisation circulaire du coefficient
commun déterminé par la preuve radiale. La troisième montre
que la réciprocité d'action et de gravitation conserve l'unité
d'aire. Ces conditions thermiques et géométriques permettent de relier
les résultats sans identifier la détermination d'une section
dimensionnelle au choix de sa coordonnée numérique.

La distinction exige de considérer la nature de chaque objet. Un rapport
sans dimension exprime une relation entre grandeurs. Une constante
dimensionnée occupe une droite de quantités et acquiert une coordonnée numérique
lors du choix d'une unité. Une condition de frontière fixe une réalisation au sein
d'une collection de solutions. Les trois déterminations peuvent intervenir
dans une même formule et conservent des fonctions différentes dans sa démonstration.
Le débat de Duff, Okun et Veneziano sur les constantes dimensionnées
illustre l'importance d'expliciter ces choix
\cite{duffokunveneziano}. La formulation actuelle du Système international
précise, quant à elle, la fonction des constantes définitoires dans la
réalisation des unités \cite{bipm2018,bipmbrochure}.

\subsection{Support arithmétique, orientation et dynamique génératrice}

Dans l'Holographie Modulaire Triadique, HMT, le support arithmétique, l'orientation
ternaire et la dynamique avec mémoire précèdent les lectures scalaires.
APP, \emph{All Possible Paths}, utilise le support discret
\(X_9=(\mathbb Z/9\mathbb Z)^2\), avec le graphe de Cayley additif
engendré par \(\pm(1,0)\) et \(\pm(0,1)\). Sur chaque position
de sa coordinatisation \(1,\ldots,9\) agissent les deux évaluations, somme et
produit, en conservant leurs restes et quotients. TRIT assigne une signature
\(\{-1,0,+1\}\) qui distingue l'orientation et le régime local.
TPK, \emph{Topological Phase Kernel}, compose sélection, transport,
mise à jour de la retenue et de la mémoire, et prolongement compatible de l'état.
Le retour de phase conserve l'avancée de la mémoire. Les opérations et la
génération régionale de \ref{sec:nucleo}--\ref{sec:eta} construisent les
régions numériques, le registre dodécaphasique et alpha avant leur utilisation dans
les lecteurs d'action, d'incidence et d'information.
La représentation analytique corrélée d'alpha prolonge à toute
profondeur la coordonnée produite par la clôture dodécaphasique.
La réalisation archimédienne détermine les valeurs de cylindres compatibles ;
la réalisation d'incidence conserve d'autres informations de la même origine.
Cette dépendance commune permet d'étudier chaque constante comme valeur et
comme structure, puis de situer la mesure comme comparaison.

Cette distinction se matérialise dès la constitution du registre.
L'évaluation des visites et traces orientées de
\ref{subsec:registro-evaluacion-incidencial} produit trois dénombrements entiers
par fenêtre, conserve leurs restes et quotients et détermine les vingt-cinq
coordonnées de \eqref{eq:registro-composicion-incidencial}. Le domaine
comprend la collection d'incidences, leurs multiplicités et leur orientation.
La proposition~\ref{prop:registro-36-residuos} caractérise exactement
l'information résiduelle qui détermine cette image ; les transformations
de \ref{subsec:k-pi-h} et \ref{subsec:k-hadamard} permettent ensuite de récupérer
le registre signé. L'évaluation du livre et la
reconstruction de sa sortie sont ainsi reliées, avec une fonction démonstrative propre à
chaque opération.

Cette structure intervient effectivement dans les équations.
Les cardinaux d'incidence interviennent dans
une fonctionnelle de retour ; les sections d'action interviennent dans une
réciprocité gravitationnelle ; les occupations de bord interviennent dans l'aire
et dans le spectre modulaire ; l'inversion d'une forme quadratique intervient
dans l'échange d'impulsion et d'enroulement. Chaque connexion s'exprime
par une application, une identité et un domaine de validité.

\subsection{Action, incidence et information de frontière}

L'argument propre de l'article II se concentre sur la relation entre
action, aire, entropie et gravité. Le paramètre de Barbero--Immirzi occupe
une place centrale parce qu'il relie le fonctionnel d'incidence à une
normalisation aréale et à sa réalisation dans l'action gravitationnelle.
La formulation de Holst fournit un langage de réalisation ultérieure
pour ce coefficient \cite{holst1996}. Dans la construction étudiée ici,
le coefficient est évalué à partir de l'incidence orientée et des canaux
angulaires déjà obtenus.

L'inclusion d'une hexade marquée dans une octade conserve la paire
d'incidence déterminant les cardinaux 90 et 120. Le tour dodécaphasique
distingue douze secteurs d'une frontière orientée de retour. Le lecteur
linéaire de cette chaîne produit les multiplicités du fonctionnel. Le
coefficient de pôle et l'évaluation décimale sont calculés ensuite.
La séquence révèle pourquoi la valeur de Barbero--Immirzi appartient à une
structure d'orientation et d'aire, en plus d'occuper une place dans
un tableau de constantes.

L'aire et le Hamiltonien modulaire décrivent des aspects complémentaires
du bord. La première lecture agrège les occupations. La seconde pondère
leur provenance sectorielle. Pour \(\Delta_{\rm mod}\ne0\) et une
normalisation connue, le Hamiltonien permet d'obtenir l'observable
pondérée \(D\) ; avec le total \(N\), celle-ci reconstruit les deux
occupations à l'origine de l'agrégat. L'inversion de la paire
\((N,D)\) n'exige pas cette condition supplémentaire. Cette récupérabilité précise
le contenu de la thèse généalogique : deux configurations peuvent
partager une lecture scalaire et conserver des différences accessibles à une autre
lecture du même état. L'argument est développé d'abord dans
l'espace fini correspondant, puis dans la réalisation déclarée.

La conservation de la mémoire précise aussi comment réutiliser une dynamique
lorsque des variables auxiliaires sont éliminées. L'identité
\eqref{mem:identidad-finita} retient l'état terminal de chaque troncature ;
la proposition~\ref{mem:prop-serie} fournit le prolongement et sa majoration
du reste. La proposition~\ref{mem:prop-schur} démontre que l'observation
complète est conservée en transformant conjointement l'opérateur réduit,
la source et le lecteur. Cette règle identifie le contenu opératoire
devant être transmis entre deux descriptions d'un même transport.

L'entropie d'intrication exige en outre une bipartition et un
état. L'exposé spécifie l'une et l'autre par une purification et son
état réduit. L'information adimensionnelle du spectre est distinguée
de l'entropie dimensionnelle, dont la normalisation incorpore la constante
de Boltzmann. Cette précision permet de relier l'énergie de la
périodicité tritique à l'action sans confondre un dénombrement d'états,
une distribution probabiliste et une température.

\subsection{Capacité, mémoire énergétique et réalisation thermométrique}

La section~\ref{sec:ii-ter27-antecedentes} développe le mécanisme fini
qui relie renouvellement des deux feuillets APP, capacité et première vacance,
double graduation du porteur et énergie de mémoire. La boucle bilatérale
détermine un spectre décimal ; son regroupement tritique conserve le reste.
La marque de capacité et les secteurs d'incidence produisent ensuite
une intensité complète dont l'origine énergétique est enregistrée. Les courants
connectés fixent les sauts entre secteurs, et le transport centré
conserve le lecteur entre réalisations d'échelles distinctes.

La section~\ref{sec:ii-termica-marcada} compose ces résultats avec l'action, l'aire, Barbero et l'horizon. Elle ajoute une réalisation thermométrique qui conserve la référence spectrale dans un registre tensoriel : le canal entropique lit l'information marquée et le canal énergétique l'espérance conditionnée. La démonstration obtient le coefficient à partir de leur quotient différentiel et démontre le principe variationnel correspondant. La référence fixe et le couple de fonctionnelles marquée--conditionnée sont les conditions constitutives de la réalisation. L'hamiltonien chronologique de capacité conserve les degrés $23,26,29$ et détermine son état normalisé ; l'application thermique de l'intensité se relie ainsi à l'horloge et à l'horizon, en maintenant leurs échelles et leurs transports effectifs.

\subsection{Composition longitudinale et détermination du couplage gravitationnel}

La gravitation fait partie de l'argument central. La composition de l'inscription, de la longueur et de la métrique fixe d'abord \(L\) et la réponse radiale. L'horloge conjuguée satisfait \(\omega_{108}L=c_{\rm int}\) ; sa période se répartit entre les 108 pas du calendrier dodécaphasique. De là proviennent \(T_{108}=2\pi_{\HMT}L/c_{\rm int}\) et \(t_0=T_{108}/108\). La condition de rayon réduit R8 détermine \(G\) sur le même caractère qui exprime déjà l'énergie. La coïncidence circulaire conserve la forme équivalente \(\vartheta_{108}^{\circ}=(54/\pi_{\HMT})^2\). Dans les sections d'action résultantes, le coefficient gravitationnel se transforme réciproquement et le produit de l'action et de la gravitation demeure invariant. L'unité d'aire est \(L^2\) dans les deux orientations.

L'égalité conservée possède une conséquence géométrique : le passage entre les sections affecte certaines coordonnées tout en préservant une quantité géométrique commune. La signification de la valeur s'obtient à partir du transport qui la produit. La relation avec l'horizon introduit ensuite l'énergie, la température et l'entropie d'aire sous des normalisations explicites. La convergence de ces quantités dans une cellule d'action se démontre par substitution des constructions précédentes.

La dépendance thermique est présentée avec la même précision. Une identité portant sur le produit de la constante de Boltzmann par une température détermine une échelle énergétique. La détermination de chaque coordonnée individuelle exige sa coordinatisation correspondante. La normalisation de chaque réalisation fixe les bases et le transport dimensionnel de cette coordonnée. La comparaison physique porte sur ces résultats et leurs unités.

\subsection{Réciprocité, propagation angulaire et continuité exceptionnelle}

Les coordonnées angulaires interviennent dans deux réalisations distinctes
de la même construction. Le fonctionnel d'incidence orientée détermine
le paramètre de Barbero--Immirzi ; le lecteur d'espèce détermine les
angles de mélange CKM et leur phase. La section~\ref{sec:ii-propagacion-angular}
expose la matrice rationnelle, son inverse et la réalisation unitaire complexe.
Le déterminant de la transformation et l'invariant de Jarlskog
fournissent des contrôles complémentaires : le premier examine la récupération
des coordonnées ; le second examine la phase des amplitudes. Les
valeurs numériques sont confrontées à l'expérience après ces constructions.

Le théorème~\ref{thm:ii-witt-angular-transport} réalise les coordonnées
angulaires dans le module d'incidence de Witt par une isométrie
explicite. Le corollaire~\ref{cor:ii-witt-composed-reader} conserve la
métrique du système sectoriel et fournit une inverse sur son image.
Le passage à l'incidence exceptionnelle transmet donc les coordonnées
produites par les règles angulaires, avec leurs relations métriques.
L'évaluation sectorielle et le transport isométrique se composent dans cet
ordre, conformément à \eqref{eq:ii-witt-composed-reader}.

L'ellipse d'action introduit une forme de déterminant fixé et un paramètre d'anisotropie. L'échange de ses demi-axes inverse le rayon modulaire. Deux transports peuvent conserver la même forme d'énergie et agir avec des signes différents sur la forme symplectique. L'orientation de l'action distingue alors des transports symplectiques et antisymplectiques ayant la même énergie scalaire. Cette propriété fonde l'attention portée au transport et à la mémoire dans la continuation du modèle.

La relation action--covariance précise le contenu quantique de cette géométrie. La proposition~\ref{prop:ii-covarianza-gaussiana} construit, dans la représentation canonique et avec l'état gaussien qui y sont spécifiés, une covariance de déterminant \(\hbar^2/4\) qui atteint l'égalité de Robertson--Schrödinger. La proposition~\ref{prop:ii-area-covarianza-reciprocidad} identifie le contour d'action à l'ellipse à deux écarts-types : ses demi-axes sont \(2\Delta Q\) et \(2\Delta P\), et son aire est \(h=2\pi\hbar\). La réciprocité échange les dispersions et conserve ce déterminant. Le passage à la position et à la quantité de mouvement maintient explicite la coordinatisation dimensionnelle, suivant~\eqref{eq:ii-covarianza-dimensiones}. L'action précédemment construite, sa distribution angulaire et les fluctuations d'une réalisation opératorielle définie sont ainsi mises en relation.

La section~\ref{delta22:ii:transporte} démontre la conservation exacte de l'action quadratique sous transport entre ellipses, dérive son terme de connexion hamiltonienne et établit la condition de positivité du générateur. La version conformément symplectique conserve l'action normalisée et explicite la forme transportée.


L'échange moment--enroulement prolonge cette réciprocité dans
un espace réticulaire muni d'un rayon. Sa démonstration conserve le quotient
de retenue, le domaine de l'opérateur et sa transformation unitaire.
La fonction modulaire correspondante est évaluée sur le rayon ainsi
transporté. L'incidence exceptionnelle est développée jusqu'à Moonshine
comme partie de la seconde réalisation de la structure commune. Les
connexions entre ces réalisations sont présentées par leurs applications
concrètes, permettant de distinguer ce qui change et ce qui se conserve.

\Needspace{7\baselineskip}
La comparaison des secteurs des quarks et des leptons exige de suivre séparément les antécédents que le développement angulaire et les opérateurs d'espèce transmettent à CKM et PMNS. Chaque transport conserve son cadre propre ; cette distinction permet de préciser où interviennent la conjugaison des amplitudes, l'orientation des parcours et les invariants de mélange. Le manuscrit maintient les différences de construction lors du passage des coordonnées structurelles aux réalisations respectives.

\subsection{Identités d'aire, d'entropie et d'échelle gravitationnelle}

La relation entre information et gravitation s'exprime par trois déterminations articulées. L'unité d'aire satisfait $\ell_P^2=G\hbar/c^3$. L'incidence orientée organise sa résolution élémentaire au moyen de $\gamma_{\HMT}a_0$. L'état réduit et la bipartition déterminent l'information accessible, dont l'expression dimensionnelle utilise $k_B$. Dans la collection sectorielle développée dans \ref{sec:ii-ecuacion-informacion}, ces dépendances prennent la forme
\begin{equation}
 \langle\widehat{\mathcal A}\rangle
 =\gamma_{\HMT}a_0\frac{G\hbar}{c^3}\langle N\rangle,
 \qquad
 S=k_B\bigl(\log Z+\Delta_{\rm mod}\langle D\rangle\bigr).
 \label{eq:ii-intro-area-informacion}
\end{equation}
Les quantités $N$ et $D$ sont les deux lectures des occupations de
bord ; $Z$ normalise l'état. Leur application conjointe récupère ces
occupations par~\eqref{eq:recover}. Les preuves adjacentes montrent
aussi comment leurs fluctuations gouvernent les réponses de l'aire et de
l'entropie.

L'étape suivante conserve ses données géométriques explicites. Dans la
réalisation d'horizon de la section~\ref{sec:confluencia},
$S_H/k_B=A_H/(4\ell_P^2)$ et $T_HS_H=E_\Lambda$.
L'aire microscopique, l'entropie de l'état réduit et l'entropie
d'horizon maintiennent leurs définitions et leurs applications de réalisation.
Cette organisation permet de lire en continuité l'incidence, l'action,
l'information et la gravitation, en précisant à chaque passage l'objet
conservé et l'observable qui change.

La section scalaire gravitationnelle possède également une réalisation tensorielle. La descente d'incidence et le projecteur de rang cinq de \ref{sec:ii-pentadica} déterminent un support à cinq composantes ; son opérateur d'entrelacement le réalise sous forme de tenseurs symétriques sans trace avant la réduction transverse de rang deux. L'orientation de ses composantes et l'observation réduite restent ainsi distinctes du couplage commun. L'aire microscopique et l'information sectorielle, construites dans \ref{sec:area} et \ref{sec:ii-ecuacion-informacion}, conservent leurs définitions lors de la composition de ces transports. La continuité exceptionnelle de \ref{exc:doble-lectura} et \ref{exc:leech} conserve l'incidence jusqu'au Moonshine ; la dualité de \ref{sec:ii-dualidad-t} conserve l'action unitaire et les domaines de l'opérateur sous l'échange de la quantité de mouvement et de l'enroulement.
% END THERMAL_R03_SYNC_01
